% fig_fo1_3_revheat_soc — 가역 발열의 SOC 부호 교대 (Ch2 §2.7·§2.8, eq:qrev·tab:qrev)
% 좌표 = 전하보존 음함수(eq:implicit) 해 U_oc(x̄)→완전식 ∂U_oc/∂T(eq:complete) 의 식 그대로 수치 평가
%   (Anode_Fit_v1.0.20.py: solve_U_oc·entropy_coefficient_x, 298.15 K). 우축 Q̇_rev/I=-T·∂U_oc/∂T.
% 대상 라이브러리(ch2): calc,arrows.meta,decorations.pathreplacing,positioning.
\begin{figure}[t]
\centering
\begin{tikzpicture}[x=8.2cm,y=4.3cm,>=Latex,font=\scriptsize]
% ---- endo 영역 옅은 음영 (x̄>x_zero) ----
\fill[black!6] (0.681,-0.40) rectangle (0.98,0.35);
% ---- 영점선 ----
\draw[densely dashed,black!55] (0.02,0) -- (0.98,0);
% ---- 축 ----
\draw[->] (0.02,-0.40) -- (0.98,-0.40) node[right] {$\bar x$};
\draw[->] (0.02,-0.40) -- (0.02,0.36) node[above,align=center,font=\tiny] {$\partial U_\oc/\partial T$\\{[mV/K]}};
\draw[->] (0.98,-0.40) -- (0.98,0.36) node[above,align=center,font=\tiny] {$\dot Q_\rev/I$\\{[mV]}};
\foreach \xx in {0.1,0.25,0.5,0.75,0.9} {\draw (\xx,-0.39)--(\xx,-0.41) node[below,font=\tiny] {\xx};}
% 좌축 눈금
\foreach \yy in {-0.3,-0.2,-0.1,0.1,0.2,0.3} {\draw (0.02,\yy)--(0.007,\yy) node[left,font=\tiny] {\yy};}
% 우축 눈금 (Q̇/I; y=∂U/∂T 에서 Q̇=-298.15·∂U/∂T)
\foreach \qq/\yy in {90/-0.3019,60/-0.2013,30/-0.1006,0/0,-30/0.1006,-60/0.2013,-90/0.3019}
  {\draw (0.98,\yy)--(0.993,\yy) node[right,font=\tiny] {$\qq$};}
% ---- config 기여 (파선) ----
\draw[densely dashed,black!60] plot[smooth] coordinates {(0.0500,-0.2273) (0.0900,-0.1730) (0.1300,-0.1377) (0.1700,-0.1108) (0.2100,-0.0888) (0.2500,-0.0700) (0.2900,-0.0533) (0.3300,-0.0382) (0.3700,-0.0243) (0.4100,-0.0113) (0.4500,0.0012) (0.4900,0.0132) (0.5300,0.0250) (0.5700,0.0369) (0.6100,0.0491) (0.6500,0.0617) (0.6900,0.0749) (0.7300,0.0887) (0.7700,0.1026) (0.8100,0.1153) (0.8500,0.1232) (0.8900,0.1223) (0.9300,0.1244) (0.9500,0.1389)};
\node[anchor=west,font=\tiny,black!60] at (0.70,0.055) {config 몫};
% ---- 완전식 ∂U_oc/∂T (굵은 실선) ----
\draw[very thick] plot[smooth] coordinates {(0.0500,-0.3737) (0.0900,-0.3175) (0.1300,-0.2799) (0.1700,-0.2506) (0.2100,-0.2259) (0.2500,-0.2039) (0.2900,-0.1838) (0.3300,-0.1649) (0.3700,-0.1466) (0.4100,-0.1289) (0.4500,-0.1113) (0.4900,-0.0936) (0.5300,-0.0756) (0.5700,-0.0570) (0.6100,-0.0376) (0.6500,-0.0170) (0.6900,0.0055) (0.7300,0.0305) (0.7700,0.0595) (0.8100,0.0946) (0.8500,0.1397) (0.8900,0.1999) (0.9100,0.2369) (0.9300,0.2790) (0.9500,0.3277)};
% ---- 5 표 앵커점 (tab:qrev) + Q̇/I 값 ----
\foreach \xx/\yy/\lab in {0.10/-0.307/{$+91.5$}, 0.25/-0.204/{$+60.8$}, 0.50/-0.089/{$+26.6$}, 0.75/0.044/{$-13.2$}, 0.90/0.218/{$-64.9$}}
  {\fill (\xx,\yy) circle (1.5pt);}
\node[anchor=south west,font=\tiny] at (0.105,-0.307) {$+91.5$};
\node[anchor=south west,font=\tiny] at (0.255,-0.204) {$+60.8$};
\node[anchor=south west,font=\tiny] at (0.505,-0.089) {$+26.6$};
\node[anchor=north east,font=\tiny] at (0.745,0.044) {$-13.2$};
\node[anchor=north east,font=\tiny] at (0.895,0.218) {$-64.9$};
% ---- 영점 교차 ----
\draw[densely dotted,black!55] (0.681,0)--(0.681,-0.40);
\node[anchor=north,font=\tiny] at (0.681,-0.008) {$\bar x_0\!\approx\!0.68$};
% ---- 발열/흡열 라벨 ----
\node[font=\tiny,align=center] at (0.30,0.20) {방전 발열 (exo)\\$\dot Q_\rev\!>\!0$\ ($\partial U_\oc/\partial T\!<\!0$)};
\node[font=\tiny,align=center] at (0.83,-0.28) {방전 흡열 (endo)\\$\dot Q_\rev\!<\!0$};
\end{tikzpicture}
\caption{탈리튬화 진행 $\bar x$ 에 따른 가역 발열의 부호 교대(발열 $\to$ 흡열) --- 좌표는 전하 보존
음함수(식~\eqref{eq:implicit})의 해 $U_\oc(\bar x)$ 를 완전식 $\partial U_\oc/\partial T$(식~\eqref{eq:complete})에
되먹인 식 그대로의 수치 평가다($298.15$ K, Chapter 1 4-전이 staging, $w_j{=}RT/F$). 굵은 실선이 완전식
$\partial U_\oc/\partial T$, 파선이 그 봉우리 내부 config 몫이며, 우축 $\dot Q_\rev/I=-T\,\partial U_\oc/\partial T$
(식~\eqref{eq:qrev})는 좌축을 $-298.15$ 배 한 것이다. 점은 표~\ref{tab:qrev} 의 다섯 앵커($\dot Q_\rev/I{=}
{+}91.5/{+}60.8/{+}26.6/{-}13.2/{-}64.9$ mV)다. 저-$\bar x$(Li-rich, $2\!\to\!1$ 큰 음의 $\Delta S$ 지배)는
$\partial U_\oc/\partial T{<}0$ 이라 방전이 발열이고, 고-$\bar x$(Li-poor)에서 config 양의 발산이 드러나 $\bar x_0
\!\approx\!0.68$ 에서 부호가 뒤집혀 흡열로 돈다 --- 임의 스위치 없이 한 종합식이 분포만으로 부호를 교대한다.
(경계 행 $\bar x\!\approx\!0.75$ 의 부호는 두-상 폭 서식 전제에 걸린 한정이 있다: 표~\ref{tab:qrev} 캡션.)}
\label{fig:cand-fo1-3}
\end{figure}
